%% Manuscript for the Journal of Statistics and Data Science Education.
%%
%% JSDSE requires the standard "article" class with no special macros, 1-inch
%% margins, double-spaced text, and natbib with plain.bst or apalike.bst.
%% Compile with pdfLaTeX. See README.md in this folder.
%%
%% Every unfinished item below is tagged "% TODO:". Before submitting, run
%%     grep -n "TODO" article.tex
%% and make sure nothing is left.

\documentclass[12pt]{article}

\usepackage[margin=1in]{geometry}
\usepackage{setspace}
\usepackage[utf8]{inputenc}
\usepackage[T1]{fontenc}
\usepackage{natbib}
\usepackage{graphicx}
\usepackage{url}

\bibliographystyle{apalike}
\doublespacing

%% ---------------------------------------------------------------------------
%% NON-ANONYMOUS VERSION. For the anonymous version required by JSDSE:
%%   - delete \author and the affiliation,
%%   - replace the repository, DOI and demo URLs with anonymised links, since
%%     they identify the author,
%%   - check that no self-citation reveals authorship.
%% Both versions must be submitted together.
%% ---------------------------------------------------------------------------

\title{\textbf{PCAPro: A Browser-Based Platform for Teaching Principal
Component Analysis with Embedded Methodological Guidance}}

\author{Luis Ángel Barrera-Guzmán\\
Centro Académico Regional sede Huatusco-Veracruz\\
Universidad Autónoma Chapingo, Huatusco, Veracruz, Mexico\\
\texttt{lbarrerag@chapingo.mx}}

\date{}

\begin{document}

\maketitle

%% ---------------------------------------------------------------------------
%% JSDSE: unstructured abstract of 200 words. Currently 193; recount if edited.
\begin{abstract}
Principal component analysis is rarely misapplied at the eigendecomposition.
It goes wrong at the decisions surrounding it: how missing values are handled,
whether variables are standardised, whether the correlation structure justifies
the analysis at all, how many components to retain, and which loadings may be
interpreted. General-purpose statistical software computes the quantities needed
to make those decisions and then stops, which is reasonable for a statistician
and is where students and non-specialist researchers go wrong. We describe
PCAPro, a platform that places the methodological argument next to the number
that justifies it. It reports eight retention criteria side by side rather than
a single answer, states which assumption is at risk and what to do about it, and
declines to present a decision as more certain than it is. PCAPro runs entirely
in a web browser with no installation, which makes it usable in shared computer
laboratories where students cannot install software, and its interface, embedded
guidance and generated report exist in full in Spanish as well as English. We
describe the pedagogical design, the implementation and its verification against
established software, and discuss the limitations of an approach that guides
without assessing.
\end{abstract}

\noindent\textbf{Keywords:} principal component analysis; statistics education;
multivariate statistics; educational software; data literacy; Spanish-language
resources

%% ---------------------------------------------------------------------------
\section{Introduction}

Principal component analysis \citep{pearson1901, hotelling1933, jolliffe2002}
is among the most widely used multivariate techniques in the biological and
agricultural sciences, and it is often the first multivariate method a student
meets. It is also among the most frequently misapplied.

The computation is not the difficulty. Every general-purpose statistical
environment provides a reliable eigendecomposition, and the numerical problem
has been settled for decades. What decides whether a principal component
analysis is defensible is a sequence of choices made before and after that
computation: how missing values are treated, whether variables are standardised,
whether the correlation structure justifies the analysis at all, how many
components are retained, whether the solution is rotated and by which method,
and which loadings are large enough to interpret. Each of these changes the
result, and none of them is decided by a function call.

Existing software computes the quantities needed to make those decisions and
then stops. FactoMineR \citep{le2008} reports the eigenvalues; deciding how many
to keep is left to the analyst. The psych package \citep{psych} reports the
Kaiser--Meyer--Olkin index; deciding whether a value of 0.62 is acceptable is
left to the analyst. This division of labour is entirely appropriate for a
statistician, who brings the missing judgement. It is precisely where students
and researchers from other disciplines go wrong, and it is invisible to them:
the software returns a result either way.

%% TODO: literature review on teaching PCA and on educational software for
%% statistics. JSDSE is an education journal and expects the work to be situated
%% in that literature, not only in the literature of the method. Without this
%% the manuscript reads as a user manual, and reviewers will say so.

Two further constraints shaped the design and are not incidental to it. First,
the software runs with no installation: it is a directory of static files served
by any web server or opened directly from disk. This is a pedagogical
requirement rather than a convenience, because a substantial part of the
intended audience works in shared university computer laboratories where nothing
can be installed and where students have no administrative rights. Second, the
interface, the embedded guidance and the generated report exist in full in
Spanish as well as English, because teaching material for multivariate
statistics in Spanish is scarce and the language barrier compounds the tooling
barrier for Latin American students.

%% ---------------------------------------------------------------------------
\section{Materials and Methods}

\subsection{Pedagogical design}

%% TODO: this is the core section for JSDSE. It must argue teaching design
%% decisions and their rationale, not describe functions. Three principles are
%% already implemented and need to be argued here:
%%   1. Guidance lives next to the computation, not in a separate manual.
%%      Thirty-four expandable sections cover the postulates, the preliminary
%%      tests, sample-size rules, the retention criteria, the rotations, the
%%      reading of cos-squared and contributions, and the three kinds of ellipse.
%%   2. Eight retention criteria are shown at once instead of a single answer,
%%      because they disagree and the disagreement is informative.
%%   3. Diagnostics do not return a bare number: they state which assumption is
%%      at risk and what to do, and refuse to present a decision as more certain
%%      than it is.

\subsection{The analysis workflow}

%% TODO: the six blocks and the decision taken at each one.

\subsection{Implementation}

PCAPro is a directory of static files that runs entirely in the browser, with
no server-side computation and no installation. All computation happens in the
user's own machine and no data are transmitted, which matters for unpublished
research data and for institutional data-protection rules.

The numerical engine was written from scratch rather than wrapped around an
existing library, because no browser-side library offered the combination of
retention criteria and rotation algorithms required. The eigendecomposition uses
cyclic Jacobi rotations. Eight retention criteria are implemented, including
Horn's parallel analysis \citep{horn1965} with the 95th-percentile rule
\citep{glorfeld1995}, Velicer's minimum average partial \citep{velicer1976},
the broken stick \citep{frontier1976} and the scree elbow \citep{cattell1966}.
Seven rotations \citep{kaiser1958, hendrickson1964} are obtained through the
gradient projection algorithm of \citet{jennrich2001, jennrich2002}. Sampling
adequacy is assessed with the Kaiser--Meyer--Olkin index \citep{kaiser1974} and
Bartlett's test of sphericity \citep{bartlett1950}.

\subsection{Verification}

Correctness is checked by an automated suite of 60 tests. Every expectation is
either a reference value published by an independent implementation or an
algebraic invariant that must hold regardless of the data. No expectation was
recorded from a previous run of PCAPro itself: comparing software against its
own past output establishes stability, not correctness. The suite runs
automatically on every change to the software, and a separate replication
script reproduces the reported results outside the browser and compares them
against R.

%% ---------------------------------------------------------------------------
\section{Results}

%% TODO: the most delicate section of this submission. Without a learning
%% evaluation, "results" cannot be a study. What is defensible:
%%   - what the software does at each decision point, illustrated with screenshots;
%%   - numerical agreement with R (table from the replication script);
%%   - availability: online demo, open source, DOI, teaching materials.
%% Worth asking the editor whether this fits as "Results" or whether the paper
%% belongs as a Brief Communication.

%% ---------------------------------------------------------------------------
\section{Discussion}

%% TODO: must state frankly:
%%   - PCAPro does not replace R and does not try to;
%%   - guiding is not assessing: whether it improves learning has not been
%%     measured, and saying so before a reviewer does;
%%   - the pedagogical risk of a tool that gives recommendations, namely that it
%%     may replace judgement instead of building it, and how showing the
%%     disagreement between criteria rather than a single answer tries to
%%     mitigate that.

%% ---------------------------------------------------------------------------
\section*{Acknowledgments}

%% TODO: write. Remove in the anonymous version.

\section*{Declaration of Interest Statement}

The author declares no competing interests.

%% ---------------------------------------------------------------------------
\section*{Software Availability}

PCAPro is free software distributed under the GNU General Public License,
version 3 or later. The source code, documentation and teaching materials are
available at \url{https://github.com/luisangelbg/PCAPro} and archived with a
persistent identifier at \url{https://doi.org/10.5281/zenodo.22649709}. A
working installation requiring no download is available at
\url{https://luisangelbg.github.io/PCAPro/}.

%% TODO: in the anonymous version, replace the three URLs above with anonymised
%% links (for example through anonymous.4open.science) or with the placeholder
%% "[repository, anonymised for review]". As written they identify the author.

\bibliography{ref}

\end{document}
